% ======================================================================
% Method Section: Prefix-Aware Tree Verification for Parallel Drafters
% ======================================================================
% Full method section for a NeurIPS-style paper.
% Requires: amsmath, amssymb, amsthm, algorithm, algorithmic, booktabs
% ======================================================================

\newtheorem{definition}{Definition}
\newtheorem{proposition}{Proposition}
\newtheorem{corollary}{Corollary}
\newtheorem{remark}{Remark}

\section{Method}
\label{sec:method}

We propose \emph{prefix-aware tree verification}, a principled approach
to constructing verification trees for speculative decoding with parallel
draft models.  Our method formulates tree construction as a submodular
maximization problem, yielding a $(1 - 1/e)$-approximate solution with
a simple greedy algorithm.

% ======================================================================
\subsection{Background: Speculative Decoding with Parallel Drafters}
\label{sec:background}
% ======================================================================

In standard speculative decoding~\citep{leviathan2023,chen2023}, an
autoregressive draft model generates $D$ candidate tokens sequentially,
and the target model verifies them in a single forward pass.  The
\emph{accepted length} $\tau$ is the number of consecutive tokens the
target agrees with before the first mismatch.

\emph{Parallel draft models} (e.g., discrete diffusion
drafters~\citep{dflash}) produce logits for all $D$ future positions
simultaneously in one forward pass, conditioned on context features from
the target model.  This is faster than autoregressive drafting but
introduces a key structural difference: \textbf{the draft logits at each
position are conditionally independent given the context}, rather than
being conditioned on previous draft tokens.

This independence means the draft produces a factored distribution
$\prod_{d=1}^{D} p_d(\sigma_d)$ over token sequences, where $p_d$ is
the draft's marginal at position~$d$.  A single draft forward pass
therefore provides logits for a combinatorially large space of candidate
sequences, from which we must select a small subset to verify
efficiently.

\paragraph{Why trees?}
Without tree-structured verification, a parallel drafter produces a
single sequence (e.g., by taking the argmax at each position) and the
target verifies it linearly.  Any mismatch at position~$k$ wastes the
verification of positions $k{+}1, \dots, D$.  A \emph{verification tree}
hedges against mismatches by proposing multiple alternatives at each
position.  The target model verifies the entire tree in one forward pass
using tree attention~\citep{miao2024specinfer,cai2024medusa}, and the
accepted length is the longest matching root-to-leaf prefix.

The challenge is: \textbf{which tree should we build?}  The tree must be
small (the target's verification cost scales with the number of trie
nodes) but cover the most probable token sequences (to maximize accepted
length).  We formalize this as an optimization problem.

% ======================================================================
\subsection{The Prefix-Coverage Objective}
\label{sec:objective}
% ======================================================================

\begin{definition}[Verification tree]
A \emph{verification tree} $T$ is a set of \emph{leaves}, where each
leaf $\ell = (\sigma_1, \dots, \sigma_{D_\ell})$ is a token sequence of
depth $D_\ell \leq D$.  The tree is stored as a trie: shared prefixes
are represented once, so the total number of trie nodes (the
verification cost) depends on the prefix overlap among leaves.
\end{definition}

\begin{definition}[Prefix-coverage objective]
\label{def:objective}
For a leaf $\ell = (\sigma_1, \dots, \sigma_{D_\ell})$, its
\emph{depth-$k$ prefix} is $\pi_k(\ell) = (\sigma_1, \dots, \sigma_k)$
for $k \leq D_\ell$.  For a set of leaves $T$, let
\[
  \mathrm{Prefixes}_k(T)
  = \bigl\{\, \pi_k(\ell) : \ell \in T,\; D_\ell \geq k \,\bigr\}
\]
be the set of distinct depth-$k$ prefixes covered by~$T$.  The
\emph{prefix-coverage objective} is
\begin{equation}
\label{eq:objective}
  f(T)
  \;=\; \sum_{k=1}^{D}\;\sum_{\sigma \,\in\, \mathrm{Prefixes}_k(T)}
        w(\sigma),
  \qquad\text{where}\quad
  w(\sigma)
  \;=\; \prod_{d=1}^{|\sigma|} p_d(\sigma_d).
\end{equation}
\end{definition}

\begin{remark}[Surrogate for expected accepted length]
\label{rem:interpretation}
Under a \emph{draft-induced product model} where the target token at
each depth~$d$ is drawn independently from $p_d$, we have
$f(T) = \mathbb{E}[\tau(T)]$, since
\[
  \mathbb{E}[\tau(T)]
  = \sum_{k=1}^{D} \Pr[\tau(T) \geq k]
  = \sum_{k=1}^{D}\;\sum_{\sigma \,\in\, \mathrm{Prefixes}_k(T)}
    \prod_{d=1}^{k} p_d(\sigma_d)
  = f(T).
\]
In practice, the target model is deterministic and the independence
assumption is approximate.  However, when the draft model is
well-calibrated, $f(T)$ correlates strongly with realized acceptance
length (see Section~\ref{sec:experiments}).  Critically, $f$ captures
the key insight that \textbf{shared prefixes should not be
double-counted}: if two leaves agree on their first $k$ tokens, they
contribute only one prefix at depth~$k$, not two.  This
prefix-awareness is what distinguishes our objective from simply ranking
leaves by cumulative probability.
\end{remark}

% ======================================================================
\subsection{Theoretical Properties}
\label{sec:theory}
% ======================================================================

\begin{proposition}[Monotone submodularity]
\label{prop:submodular}
The prefix-coverage objective $f \colon 2^{\mathcal{C}} \to
\mathbb{R}_{\geq 0}$ is monotone non-decreasing and submodular.
\end{proposition}

\begin{proof}
Express $f$ as a non-negative weighted coverage function.  Let
$\mathcal{U} = \bigcup_{k=1}^{D} \mathcal{V}^k$ be the universe of all
possible prefixes up to depth~$D$.  Each element $\sigma \in
\mathcal{U}$ has weight $w(\sigma) \geq 0$.  Each leaf $\ell \in
\mathcal{C}$ covers the set $S(\ell) = \{\pi_1(\ell), \dots,
\pi_{D_\ell}(\ell)\}$.  Then
\[
  f(T) = \sum_{\sigma \in \mathcal{U}} w(\sigma)
    \cdot \mathbf{1}\!\bigl[\sigma \in \textstyle\bigcup_{\ell \in T}
    S(\ell)\bigr].
\]

\emph{Monotonicity}: Adding a leaf can only enlarge the covered set, so
$f(T \cup \{\ell\}) \geq f(T)$.

\emph{Submodularity}: The marginal gain of adding $\ell$ to $T$ is
\[
  \Delta(\ell \mid T)
  = \sum_{k=1}^{D_\ell}
    \mathbf{1}\!\bigl[\pi_k(\ell) \notin \mathrm{Prefixes}_k(T)\bigr]
    \cdot w\!\bigl(\pi_k(\ell)\bigr).
\]
For $A \subseteq B$, we have $\mathrm{Prefixes}_k(A) \subseteq
\mathrm{Prefixes}_k(B)$, so every indicator that is~$1$ for~$B$ is
also~$1$ for~$A$.  Hence $\Delta(\ell \mid A) \geq \Delta(\ell \mid
B)$, the definition of submodularity.
\end{proof}

\begin{corollary}[Approximation guarantee]
\label{cor:greedy}
The greedy algorithm that iteratively selects the candidate with maximum
marginal gain, starting from $T_0 = \emptyset$, produces a set $T_M$ of
at most $M$ leaves satisfying
\[
  f(T_M) \;\geq\; \Bigl(1 - \frac{1}{e}\Bigr) \cdot
  \max_{|T| \leq M} f(T).
\]
\end{corollary}

\begin{proof}
Direct application of \citet{nemhauser1978}: $f$ is monotone
submodular with $f(\emptyset) = 0$, and we impose a cardinality
constraint $|T| \leq M$.
\end{proof}

% ======================================================================
\subsection{Algorithm: Prefix-Aware Tree Construction}
\label{sec:algorithm}
% ======================================================================

Our algorithm operates in three phases: (1)~candidate generation,
(2)~submodular leaf selection, and (3)~trie packing.

% ------ Phase 1 ------
\paragraph{Phase~1: Candidate generation.}
We construct a pool of $N \approx 2$--$3M$ candidate leaves via
best-first search over the draft probability space.  Starting from the
top-$K$ tokens at depth~1, a priority queue (ordered by cumulative
log-probability) repeatedly expands the highest-scoring partial sequence
by appending each of the top-$K$ tokens at the next depth.  Expansion
terminates when the pool reaches the budget $N = \max(2M, M + KD)$.

This phase explores the high-probability region of the exponentially
large token space $\mathcal{V}^D$ without enumerating it.  The pool
budget is deliberately larger than $M$ so that Phase~2 has sufficient
diversity to exploit prefix coverage.

% ------ Phase 2 ------
\paragraph{Phase~2: Greedy submodular selection.}
From the candidate pool $\mathcal{C}$ of $N$ leaves, we select $M$
leaves to maximize~$f(T)$ using the greedy algorithm guaranteed by
Corollary~\ref{cor:greedy}.

At each iteration, we select the candidate $\ell$ with the largest
marginal gain
\[
  \Delta(\ell \mid T_i) = \sum_{k=1}^{D_\ell}
    \mathbf{1}\bigl[\pi_k(\ell) \notin \mathrm{Prefixes}_k(T_i)\bigr]
    \cdot w\bigl(\pi_k(\ell)\bigr),
\]
add it to $T_{i+1}$, and mark all its prefixes as covered.

We accelerate this with the \emph{lazy greedy} technique
\citep{minoux1978}: a max-heap stores stale upper bounds on marginal
gains.  When a candidate is popped from the heap, its gain is
recomputed; if it is still the largest, it is accepted; otherwise it is
re-inserted with its updated gain.  By submodularity, marginal gains
can only decrease as $T_i$ grows, so stale values are valid upper
bounds.  Lazy greedy returns the \emph{identical} solution as standard
greedy while typically requiring far fewer gain evaluations.

% ------ Phase 3 ------
\paragraph{Phase~3: Trie packing.}
The $M$ selected leaves are inserted into a trie data structure.
Leaves that share a prefix of length~$k$ share the same trie nodes for
positions $1, \dots, k$.  The trie is serialized into:
\begin{itemize}[nosep,leftmargin=*]
  \item \texttt{packed\_ids}: token IDs in BFS order,
  \item \texttt{packed\_pos}: position IDs for each node,
  \item \texttt{parent\_idx}: parent pointer for each node (root has $-1$).
\end{itemize}
These are passed to the target model for a single forward pass with a
tree attention mask (Section~\ref{sec:tree-attention}).

% ------ Pseudocode ------
\begin{algorithm}[t]
\caption{Prefix-Aware Tree Construction}
\label{alg:prefix-aware}
\begin{algorithmic}[1]
\Require Draft logits $\{p_d\}_{d=1}^D$, leaf budget $M$, expansion width $K$
\Ensure Packed trie for one-pass target verification

\Statex \textbf{Phase 1: Candidate generation}
\State $\mathcal{F} \leftarrow$ priority queue seeded with top-$K$ tokens at depth~1
\State $\mathcal{C} \leftarrow \emptyset$, \; $N \leftarrow \max(2M,\; M + KD)$
\While{$|\mathcal{C}| + |\mathcal{F}| < N$}
  \State Pop highest-probability partial path $(\boldsymbol{\sigma}, d)$ from $\mathcal{F}$
  \If{$d = D$}
    \State $\mathcal{C} \leftarrow \mathcal{C} \cup \{\boldsymbol{\sigma}\}$
  \Else
    \For{each of top-$K$ tokens $v$ at depth $d{+}1$}
      \State Push $(\boldsymbol{\sigma} \| v,\; d{+}1)$ into $\mathcal{F}$
    \EndFor
  \EndIf
\EndWhile
\State Move remaining frontier items into $\mathcal{C}$

\Statex \textbf{Phase 2: Greedy submodular selection}
\State $T \leftarrow \emptyset$, \; $\mathrm{Covered}[k] \leftarrow \emptyset$ for $k = 1, \dots, D$
\State Initialize max-heap $\mathcal{H}$ with $\Delta(\ell \mid \emptyset)$ for all $\ell \in \mathcal{C}$
\While{$|T| < M$ and $\mathcal{H} \neq \emptyset$}
  \State Pop candidate $\ell$ with largest (stale) gain from $\mathcal{H}$
  \State Recompute $g \leftarrow \Delta(\ell \mid T)$
  \If{$g \geq$ top of $\mathcal{H}$} \Comment{Lazy greedy: accept if still best}
    \State $T \leftarrow T \cup \{\ell\}$
    \For{$k = 1, \dots, D_\ell$}
      \State $\mathrm{Covered}[k] \leftarrow \mathrm{Covered}[k] \cup \{\pi_k(\ell)\}$
    \EndFor
  \Else
    \State Push $\ell$ back into $\mathcal{H}$ with gain $g$
  \EndIf
\EndWhile

\Statex \textbf{Phase 3: Trie packing}
\State Insert all $\ell \in T$ into trie, merging shared prefixes
\State Serialize to \texttt{packed\_ids}, \texttt{packed\_pos}, \texttt{parent\_idx}
\State \Return packed trie tensors
\end{algorithmic}
\end{algorithm}

% ======================================================================
\subsection{Tree Attention and Verification}
\label{sec:tree-attention}
% ======================================================================

Given the packed trie from Algorithm~\ref{alg:prefix-aware}, we
construct a tree attention mask that enforces two constraints:
(1)~each node can attend to its ancestors and itself, and
(2)~nodes on different branches cannot attend to each other.  This mask
is derived from the \texttt{parent\_idx} array: node~$q$ can attend to
node~$k$ if and only if $k$ lies on the path from the root to~$q$.

The target model performs a single forward pass over the packed trie
with this attention mask, producing logits at every node.  We then
select the leaf whose root-to-leaf token sequence has the longest
prefix matching the target's greedy (or sampled) output.  The
accepted tokens are appended to the generated sequence, and the
KV-cache is surgically trimmed to retain only the accepted path.

% ======================================================================
\subsection{Computational Overhead}
\label{sec:overhead}
% ======================================================================

Phase~1 performs $O(N \log N)$ heap operations over candidate sequences
of length at most~$D$.  Phase~2 performs at most $O(NM)$ marginal-gain
evaluations, each scanning $D$ prefix depths, for a worst-case cost of
$O(NMD)$.  Lazy greedy substantially reduces this in practice.  Phase~3
(trie insertion) is $O(MD)$.

The entire tree construction runs on CPU with no GPU synchronization.
In our experiments (Section~\ref{sec:experiments}), tree construction
accounts for \textbf{1.9\%--2.1\%} of total decode time, while the
target model backbone accounts for 71\%--73\%.  Tree construction
overhead is negligible.

% ======================================================================
\subsection{Comparison with Prior Tree Construction Methods}
\label{sec:comparison}
% ======================================================================

To situate our approach, we compare with three alternative strategies
that we implemented and evaluated.  All use the same parallel draft
logits and tree attention verification pipeline; they differ only in
how the tree is constructed.

\paragraph{Threshold + Cartesian cap (v1).}
At each position, a confidence threshold determines the branching
count (1--$K$), and the tree enumerates the Cartesian product of
per-position candidates.  This requires three hand-tuned thresholds
($\theta_{\text{uni}}, \theta_{\text{bi}}, \theta_{\text{tri}}$) and
produces exponentially large tries: in our experiments, the average
trie size is 188 nodes (vs.\ 70 for our method at the same
\texttt{max\_tree\_size}).  Despite higher $\bar{\tau}$ than sequential
verification, the excessive verification cost often results in
\emph{lower} wall-clock speedup.

\paragraph{Expand + rerank (v2).}
Inspired by EAGLE-2~\citep{li2024eagle2}, this method expands the
tree layer-by-layer, selecting the top-$K$ nodes per layer by
cumulative confidence, then globally reranks all nodes and keeps the
top-$M$.  This produces compact trees (33 nodes at $M{=}32$) but does
not account for prefix overlap: two high-probability leaves that share
a long prefix are both selected, wasting coverage.

\paragraph{Best-first (v3).}
A priority queue expands the highest-probability frontier node and
stops when the leaf count reaches~$M$.  This is equivalent to our
Phase~1 with the pool budget set to exactly~$M$.  It produces
reasonable trees but optimizes cumulative probability rather than
prefix coverage, missing the submodular structure.

\paragraph{Prefix-aware (ours).}
By decoupling candidate generation (Phase~1) from selection
(Phase~2), our method generates a diverse pool and then selects
leaves that \emph{maximize prefix coverage}, explicitly accounting for
shared prefixes.  The submodular formulation provides a
$(1 - 1/e)$ approximation guarantee that none of the above methods
enjoy.  Table~\ref{tab:method-comparison} summarizes the key
differences.

\begin{table}[t]
\centering
\small
\caption{Comparison of tree construction strategies.  Our method (v4)
is the only one with a formal approximation guarantee and requires the
fewest human-tuned parameters.}
\label{tab:method-comparison}
\begin{tabular}{@{}lccc@{}}
\toprule
Method & Parameters & Objective & Guarantee \\
\midrule
Threshold (v1)    & 5 ($\theta$s, $K$, $M$) & Heuristic       & None \\
Expand+rerank (v2)& 2 ($K$, $M$)            & Cum.\ prob.\    & None \\
Best-first (v3)   & 2 ($K$, $M$)            & Cum.\ prob.\    & None \\
Prefix-aware (v4) & 2 ($K$, $M$)            & Prefix coverage & $(1{-}1/e)$-approx \\
\bottomrule
\end{tabular}
\end{table}

% ======================================================================
% References needed (add to main .bib):
% ======================================================================
% nemhauser1978    — submodular maximization
% minoux1978       — lazy greedy acceleration
% leviathan2023    — speculative decoding
% chen2023         — speculative sampling
% li2024eagle2     — EAGLE-2
% dflash           — DFlash (your parallel drafter)
% miao2024specinfer — SpecInfer (tree attention)
% cai2024medusa    — Medusa (tree verification)
